\documentclass[12pt]{article}

\usepackage[a4paper,margin=1in]{geometry}
\usepackage{amsmath,amsfonts,amssymb}
\usepackage{graphicx}
\usepackage{hyperref}
\usepackage{enumitem}

\title{\textbf{Lecture 15 Scribe}\\
CSE400 -- Fundamentals of Probability in Computing\\
\large Joint Random Variables and Joint Distributions}
\author{}
\date{}

\begin{document}

\maketitle

\section*{1. Why Study Joint Random Variables?}

In science and in real life, we are often interested in \textbf{two (or more) random variables at the same time}.

Examples include:

\begin{itemize}
\item Height and weight of giraffes
\item IQ and birthweight of children
\item Frequency of exercise and rate of heart disease
\item Air pollution level and respiratory illness rate
\item Number of Facebook friends and age of members
\end{itemize}

These situations involve \textbf{relationships between variables}, which require joint probability models.

\subsection*{Example 1: Smart Transportation}

Two random variables:

\begin{itemize}
\item $X =$ Vehicle Speed
\item $Y =$ Traffic Density
\end{itemize}

Question: Can speed be high when the traffic density is also high?

These variables are \textbf{correlated random variables}.

Applications in intelligent transportation systems include:

\begin{itemize}
\item Adaptive traffic signals
\item Congestion prediction
\item Autonomous vehicle routing
\end{itemize}

\subsection*{Example 2: Wireless Network Performance}

Two random variables:

\begin{itemize}
\item $X =$ Signal-to-Noise Ratio (SNR)
\item $Y =$ Data Throughput
\end{itemize}

Throughput depends on SNR, but network congestion may affect this relationship.

Applications include:

\begin{itemize}
\item 5G scheduling
\item Adaptive modulation
\item Network planning
\end{itemize}

\subsection*{Additional Examples}

\begin{itemize}
\item Weather prediction: $X =$ Rainfall, $Y =$ Temperature
\item Drone delivery systems: $X =$ Wind speed, $Y =$ Delivery time
\item Rolling two dice: events like $(X+Y=7)$ and $(X>Y)$ require a joint PMF
\end{itemize}

\subsection*{Importance of Joint Distributions}

When multiple random variables are involved:

\begin{itemize}
\item A \textbf{joint distribution} allows computation of probabilities involving both variables.
\item It helps understand \textbf{relationships between variables}.
\item If variables are independent, analysis is simpler.
\item If not independent, \textbf{covariance and correlation} measure dependence.
\end{itemize}

\section*{2. Discrete Case -- Joint PMF}

\subsection*{Setup}

Suppose $X$ and $Y$ are discrete random variables.

\[
X = \{x_1, x_2, \dots, x_n\}, \quad
Y = \{y_1, y_2, \dots, y_m\}
\]

The ordered pair $(X,Y)$ takes values in the \textbf{product set}

\[
\{(x_1,y_1),(x_1,y_2),\dots,(x_n,y_m)\}
\]

\subsection*{Definition: Joint PMF}

The \textbf{joint probability mass function (PMF)} is defined as

\[
p(x_i, y_j) = P(X = x_i, Y = y_j)
\]

This gives the probability of the \textbf{joint outcome occurring simultaneously}.

\section*{3. Joint PMF Table Representation}

The joint PMF can be organized in a table.

\begin{center}
\begin{tabular}{c|c c c}
$X \backslash Y$ & $y_1$ & $y_2$ & $y_j$ \\
\hline
$x_1$ & $p(x_1,y_1)$ & $p(x_1,y_2)$ & $p(x_1,y_j)$ \\
$x_2$ & $p(x_2,y_1)$ & $p(x_2,y_2)$ & $p(x_2,y_j)$ \\
$\vdots$ & $\vdots$ & $\vdots$ & $\vdots$ \\
$x_n$ & $p(x_n,y_1)$ & $p(x_n,y_2)$ & $p(x_n,y_j)$ \\
\end{tabular}
\end{center}

\subsection*{Key Properties}

\[
0 \le p(x_i, y_j) \le 1
\]

\[
\sum_{i=1}^{n}\sum_{j=1}^{m} p(x_i, y_j) = 1
\]

\section*{4. Example: Rolling Two Dice}

\textbf{Experiment}

Roll two fair dice.

Define:

\begin{itemize}
\item $X =$ value on the first die
\item $T =$ total (sum) of both dice
\end{itemize}

Each outcome $(i,j)$ has probability

\[
\frac{1}{36}
\]

Thus we construct the \textbf{joint probability table of $(X,T)$}.

\textbf{Observations}

\begin{itemize}
\item Entries are either $0$ or $\frac{1}{36}$.
\item $X$ and $T$ are \textbf{not independent}.
\end{itemize}

\section*{5. Continuous Case -- Joint PDF}

The continuous case is analogous to the discrete case.

Correspondence:

\begin{itemize}
\item Discrete values $\rightarrow$ Continuous intervals
\item Joint PMF $\rightarrow$ Joint PDF
\item Sums $\rightarrow$ Double integrals
\end{itemize}

If

\[
X \in [a,b], \quad Y \in [c,d]
\]

then the pair $(X,Y)$ takes values in

\[
[a,b] \times [c,d]
\]

\subsection*{Definition: Joint PDF}

A function $f(x,y)$ is called the \textbf{joint probability density function} of $X$ and $Y$ if

\[
P((X,Y) \text{ lies in a small rectangle of area } dx\,dy)
\approx f(x,y)\,dx\,dy
\]

\subsection*{Key Properties}

\[
f(x,y) \ge 0
\]

\[
\int_c^d \int_a^b f(x,y)\,dx\,dy = 1
\]

Note: The joint PDF $f(x,y)$ may take values greater than 1 because it represents a \textbf{density}, not a probability.

\section*{6. Example}

Given

\[
f(x,y) = 4xy, \quad 0 \le x \le 1,\; 0 \le y \le 1
\]

The pair $(X,Y)$ takes values in the unit square $[0,1]^2$.

\subsection*{Step 1: Check Validity}

\[
\int_0^1 \int_0^1 4xy\,dx\,dy
\]

\[
= \int_0^1 \left[\int_0^1 4xy\,dx \right] dy
\]

\[
= \int_0^1 2y\,dy
\]

\[
= [y^2]_0^1 = 1
\]

Hence

\begin{itemize}
\item $f(x,y) \ge 0$
\item Total probability = 1
\end{itemize}

Therefore it is a \textbf{valid joint PDF}.

\subsection*{Step 2: Compute $P(A)$}

\[
P(A) =
\int_0^{0.5} \int_{0.5}^{1} 4xy\,dy\,dx
\]

\[
= \int_0^{0.5} [2xy^2]_{0.5}^{1} dx
\]

\[
= \int_0^{0.5} 2x(1 - 0.25)\,dx
\]

\[
= \int_0^{0.5} \frac{3}{2}x\,dx
\]

\[
= \frac{3}{16}
\]

Thus

\[
P(A) = \frac{3}{16}
\]

\section*{7. Joint Cumulative Distribution Function (Joint CDF)}

Suppose $X$ and $Y$ are jointly distributed random variables.

The notation

\[
X \le x,\quad Y \le y
\]

represents the event

\[
\{X \le x \text{ and } Y \le y\}
\]

\subsection*{Definition}

The \textbf{joint cumulative distribution function} is

\[
F(x,y) = P(X \le x, Y \le y)
\]

\section*{8. Joint CDF -- Continuous Case}

If $X$ and $Y$ are continuous random variables with joint density $f(x,y)$ over $[a,b]\times[c,d]$, then

\[
F(x,y) = \int_c^y \int_a^x f(u,v)\,du\,dv
\]

\subsection*{Recovering the Joint PDF}

\[
f(x,y) = \frac{\partial^2}{\partial x \partial y} F(x,y)
\]

\section*{9. Joint CDF -- Discrete Case}

If $X$ and $Y$ are discrete random variables with joint PMF $p(x_i,y_j)$, then

\[
F(x,y) = \sum_{x_i \le x} \sum_{y_j \le y} p(x_i,y_j)
\]

\subsection*{Interpretation}

Add all probabilities in the table

\begin{itemize}
\item Rows with $x_i \le x$
\item Columns with $y_j \le y$
\end{itemize}

\section*{10. Properties of the Joint CDF}

\subsection*{1. Non-decreasing}

$F(x,y)$ is \textbf{non-decreasing}.  
If $x$ or $y$ increases, $F(x,y)$ must stay constant or increase.

\subsection*{2. Lower-Left Corner Property}

\[
F(x,y) = 0
\]

at the lower-left of the joint range.

If the lower-left corner is $(-\infty,-\infty)$,

\[
\lim_{(x,y)\to(-\infty,-\infty)} F(x,y) = 0
\]

\subsection*{3. Upper-Right Corner Property}

\[
F(x,y) = 1
\]

at the upper-right of the joint range.

If the upper-right corner is $(\infty,\infty)$,

\[
\lim_{(x,y)\to(\infty,\infty)} F(x,y) = 1
\]

\end{document}
